\section*{Chapter \ref{ch:s1:value-quality-and-low-risk} --- Value, Quality and Low Risk}

\begin{solution}{exo:s1:value-quality-and-low-risk:1}
The gain at the peak was \$8.77. A 57.8\% fall leaves $9.77 \times 0.422 = \$4.12$, so the loss is \$5.65 of the \$8.77 gained: 64\%.
\end{solution}

\begin{solution}{exo:s1:value-quality-and-low-risk:2}
Long $1/0.6 = \$1.67$ of low-beta stocks and short $1/1.5 = \$0.67$ of high-beta stocks: each side has a beta of one, and the book is net long
$1.67 - 0.67 = \$1.00$, financed by borrowing.
\end{solution}

\begin{solution}{exo:s1:value-quality-and-low-risk:3}
Each side earns $0.5 \times 1.75 \times 3\%$, so the book earns $2 \times 0.5 \times 1.75 \times 0.03 = 5.25\%$ a year before costs, if the premium
were the only thing its stocks carried.
\end{solution}

\begin{solution}{exo:s1:value-quality-and-low-risk:4}
Book changes slowly, so most of the month-to-month change in book-to-price is minus the stock's return: stocks that have fallen become cheap. A current-price
value book is therefore long recent losers and short recent winners, the opposite of momentum. Two strategies with positive premiums and negatively
correlated returns combine into a portfolio with a higher Sharpe ratio than either.
\end{solution}

\begin{solution}{exo:s1:value-quality-and-low-risk:5}
Its signal, earnings over book, falls as book rises, so it ranks high-book (cheap) stocks low: the book is short value, and the synthetic market pays a value
premium. Its correlation with the value book is $-0.60$ and its signal's IC against the planted value premium $-0.29$.
\end{solution}

\begin{solution}{exo:s1:value-quality-and-low-risk:6}
$0.594 \times 7.2\% = 4.3\%$ and $1.604 \times 7.2\% = 11.5\%$. The deciles earned 6.8\% and 8.6\%: 2.5 points more at low beta and 2.9 points less at high
beta than the model says. The security market line is flat.
\end{solution}

\begin{solution}{exo:s1:value-quality-and-low-risk:7}
The composite's Sharpe ratio is 0.44. The current-price value book and the momentum book correlate $-0.56$, so their signals partly cancel each other's
noise: the composite keeps both premiums with less risk than either book alone.
\end{solution}

\begin{solution}{exo:s1:value-quality-and-low-risk:8}
Fourteen years of losses are compatible with a premium that still exists: HML's standard deviation is 10\% a year, so a 3.6\% mean has a standard error of
about 2.7\% over fourteen years, and value earned 9.4\% a year from 2021. The drought is evidence to weigh, not proof; the better questions are whether its
measurement still captures value (intangibles), and how value fits the rest of the book.
\end{solution}

\begin{solution}{pb:s1:value-quality-and-low-risk:1}
\begin{enumerate}
\item Book equity over price; buying high ratios and selling low ones. HML is half small value plus half big value minus half small growth and half big growth,
rebuilt each June with the last fiscal year's book and December's market equity.
\item 3.6\% a year and 0.35 over 1963--2026; 5.7\% to 2006, $-5.3\%$ in 2007--2020, 9.4\% since 2021.
\item From \$9.77 at the end of 2006 down 57.8\% to September 2020; still 29.6\% below the peak in July 2026.
\item Book-to-price with an old price is a slow value signal; with today's price it also loads on recent returns (short momentum).
\item Long profitable, short unprofitable firms; quality adds growth and safety. Gross profitability predicts returns about as well as book-to-market; QMJ
earns significant risk-adjusted returns in 25 countries.
\item 0.09, $-0.19$ and 0.68 over 1963--2026; $-0.14$, $-0.48$ and 0.49 over 2007--2020.
\item 2.8\% a year at a Sharpe ratio of 0.51.
\item It is short value in disguise, and the synthetic market pays value but not profitability.
\item Low-risk stocks earn more risk-adjusted than the capital asset pricing model predicts; BAB is long levered low beta and short delevered high beta.
\item The lowest decile (beta 0.59) earned 6.8\% over bills and the highest (1.60) 8.6\%, against 4.3\% and 11.5\% predicted.
\item 4.3\% a year from 1968, Sharpe ratio 0.26, beta $-0.06$.
\item The synthetic market's security market line is steep (2.4\% at beta 0.49, 13.4\% at 1.39): nothing rewards low beta there.
\item \textbf{Named result.} HML fell 57.8\% from its December 2006 peak to September 2020; its correlation with RMW was 0.09 over 1963--2026 and $-0.14$ over
2007--2020.
\item Fama--French timing 0.24, yearly with today's price 0.21, look-ahead $-0.06$, monthly with today's price $-0.07$: the more the signal folds in recent
returns, the more it is short the planted momentum.
\item It reaches 0.44, because current-price value and momentum correlate $-0.56$.
\item Timely prices forecast true book-to-price better and earned alphas of 3.05--3.78 points a year against the standard method.
\item Varying factor weights on forecasts of their returns. For: valuation spreads and the price of quality have predicted returns. Against: slow signals,
short samples, and value's cheapness through 2007--2020.
\item Filing dates of the fundamentals, the timing of prices, the drought years in the sample, and correlations with the other factors held.
\item Near-fixed weights across value, momentum and profitability, each scaled by volatility, with timing at most a small tilt.
\item Each is a slow premium with long periods of losses that its partners tend to offset.
\end{enumerate}
\end{solution}

\begin{solution}{iq:s1:value-quality-and-low-risk:1}
Each June, stocks are split at the NYSE median size and at the 30th and 70th NYSE percentiles of book-to-market (last fiscal year's book over December's
market equity); HML is the average of the small and big value portfolios minus the average of the small and big growth portfolios, value-weighted, held for a
year.
\end{solution}

\begin{solution}{iq:s1:value-quality-and-low-risk:2}
Investors who cannot or will not use leverage get market exposure by buying high-beta stocks, bidding them up; benchmarked managers prefer them too. High beta
then comes with low alpha, and the security market line is flatter than the capital asset pricing model says.
\end{solution}

\begin{solution}{iq:s1:value-quality-and-low-risk:3}
That value had a long, deep drawdown, which a 10\% volatility factor can have by chance with a positive premium; that growth, low rates and intangibles
dominated the period. Not that value is dead: fourteen years are not enough to tell, and it recovered part of the loss from 2021.
\end{solution}

\begin{solution}{iq:s1:value-quality-and-low-risk:4}
Their returns are uncorrelated or negatively correlated with value's, so the combination keeps the premiums with less risk; profitability also separates cheap
good firms from cheap bad ones, which improves value.
\end{solution}

\begin{solution}{iq:s1:value-quality-and-low-risk:5}
Book values known only from their filing dates, not their fiscal ends; restatements; the price used (current or at the fiscal end); share counts and splits
between the book's date and today; delisted firms kept in the history.
\end{solution}

\begin{solution}{iq:s1:value-quality-and-low-risk:6}
With equal weights and equal volatilities $\sigma$, the mean is $0.3\sigma$ and the variance $\frac14\sigma^2(2 + 2\rho) = \frac14\sigma^2$, so the Sharpe
ratio is $0.3/0.5 = 0.6$. The best combination has $SR^2 = s^\top C^{-1} s = 2 \times 0.09/(1 + \rho) = 0.36$, again 0.6: by symmetry equal weights are
optimal.
\end{solution}
